\section{Joint trace control on the original source and relation spaces}
\label{sec:signed-projection-control}

Retain \(h,k,\chi=\chi_{h,k}\) from the marked-product note. Here
\(q=\deg\chi\geq1\) denotes the polynomial degree. Use the original
coordinate \(S=k/2+iy\), the densities \(m=w_h^{*k}\) and
\(r=r_{1/4}^{*k}\) with their literal masses, and
\(\mu_x=((1-x)r+xm)\,dy\) for \(0\leq x\leq1\).
The common source is \(H=\mathcal P_{2q}\) in the retained monomial basis.
Its positive definite Gram and its original metric derivative are
\[
 M_x=M_0+x\dot M,\quad M_0=M^\Gamma_{2q},\quad
 \dot M=M^{\rm ar}_{2q}-M^\Gamma_{2q},\quad
 T_x=M_x^{-1}\dot M .
 \tag{SP1}
\]
The equality \(M_xT_x=\dot M=\dot M^*\) proves that \(T_x\) is
self-adjoint in the \(M_x\) inner product. It is not asserted positive.

\subsection{All four endpoints as operators on one source}
Let \(I_N:\mathcal P_N\to H\) be literal coefficient inclusion, and
let \(B_N:\mathcal P_{N-q}\to H\) be multiplication by the original
monic \(\chi\), followed by coefficient inclusion. Define
\[
 P_N=I_N(I_N^*M_xI_N)^{-1}I_N^*M_x,\qquad
 Q_N=B_N(B_N^*M_xB_N)^{-1}B_N^*M_x,\quad Q_{q-1}=0 .
 \tag{SP2}
\]
The second formula is used for \(N\geq q\); its Gram is invertible
because multiplication by \(\chi\) is injective on polynomials.
Direct multiplication proves that these are the \(M_x\)-orthogonal
projections onto the original source and relation subspaces.
Within either family the subspaces are nested, so their orthogonal
projections commute and their product is the smaller projection.
Indeed the smaller subspace, its orthogonal complement inside the
larger one, and the orthogonal complement of the larger one give
three invariant summands on which this statement is immediate.

Let \(V_N(x)\) be the original canonical quotient volume. In the
fixed quotient-lift and monic-relation coordinates, block determinant
gives
\[
 \log V_N(x)=\log\det(I_N^*M_xI_N)
                  -\log\det(B_N^*M_xB_N).
 \tag{SP3}
\]
The relation determinant is one when its dimension is zero.
The fixed monic change-of-basis determinant has modulus one, so
SP3 retains every volume factor. Differentiation and cyclicity of
finite matrix trace give
\[
 (\log V_N)'=\operatorname{Tr}(T_x(P_N-Q_N)).
 \tag{SP4}
\]
For example
\(\operatorname{Tr}(T_xP_N)=
\operatorname{Tr}((I_N^*M_xI_N)^{-1}I_N^*\dot M I_N)\);
the relation formula uses precisely the same calculation.
Set
\[
 F(x)=\log\frac{V_{q-1}(x)V_q(x)}
                      {V_{2q-1}(x)V_{2q}(x)},\qquad
 \Delta_{h,k}=F(1)-F(0),
 \tag{SP5}
\]
and define two positive operators on this same source:
\[
 \mathcal R_x=Q_{2q-1}+Q_{2q}-Q_q,\qquad
 \mathcal P_x=(P_{2q-1}-P_{q-1})+(P_{2q}-P_q),\qquad
 \mathcal A_x=\mathcal R_x-\mathcal P_x .
 \tag{SP6}
\]
The four original signs in SP4 give
\[
                  F'(x)=\operatorname{Tr}(T_x\mathcal A_x).
 \tag{SP7}
\]
For the original row \(v=v_{2q}(y)\), the density identity has the exact
metric factor
\[
 \mathsf R_x(y)-\mathsf P_x(y)=v\,\mathcal A_xM_x^{-1}v^* .
\]
Indeed \(P_NM_x^{-1}=I_N(I_N^*M_xI_N)^{-1}I_N^*\),
and the same identity holds with \(B_N,Q_N\). This recovers the
marked-product kernels and proves the common-source matrix comparison.

\subsection{Ranks, multiplicities and the retained overlap}
The nested relation spaces in SP6 have dimensions \(1,q,q+1\).
On the first space \(\mathcal R_x\) has eigenvalue one; on the next
orthogonal increment it has eigenvalue two with multiplicity \(q-1\);
on the last increment it has eigenvalue one; elsewhere it is zero.
The nested source increments in \(\mathcal P_x\) give the same list.
When \(q=1\), the multiplicity \(q-1\) is zero. Consequently
\[
 \begin{gathered}
 \operatorname{rank}\mathcal R_x=\operatorname{rank}\mathcal P_x=q+1,
 \quad
 \operatorname{Tr}\mathcal R_x=\operatorname{Tr}\mathcal P_x=2q,\\
 \operatorname{Tr}\mathcal R_x^2
       =\operatorname{Tr}\mathcal P_x^2=4q-2,\quad
 \operatorname{Tr}\mathcal A_x=0,\\
 \operatorname{Tr}\mathcal A_x^2
       =8q-4-2\operatorname{Tr}(\mathcal R_x\mathcal P_x).
 \end{gathered}
 \tag{SP8}
\]
The equal kernel masses \(2q\) are traces. The two positive operators
have rank \(q+1\), with their multiplicity-two directions retained.

The overlap in SP8 is at least one. Each positive operator is at
least the orthogonal projection onto its range, since its nonzero
eigenvalues are one and two. The two ranges each have dimension
\(q+1\) in dimension \(2q+1\), so their intersection has dimension at
least one. Choose a unit vector in this intersection and complete
it to an orthonormal basis of the first range. The trace of the
product of the two range projections is the sum of squared norms
of the second projection on this basis, hence at least one.
For positive operators, replacing one factor by a larger positive
factor cannot decrease the trace of the product: the difference
is \(\operatorname{Tr}(B^{1/2}(A'-A)B^{1/2})\geq0\).
Applying this twice proves the stated bound for the actual operators.

\subsection{A bound that uses cancellation before estimating}
Write
\(d_x=\lambda_{\max}(T_x)-\lambda_{\min}(T_x)\).
Since \(\operatorname{Tr}\mathcal A_x=0\), for every real \(c\)
\[
       F'(x)=\operatorname{Tr}((T_x-cI)\mathcal A_x).
 \tag{SP9}
\]
Choose \(c\) to be the midpoint of the two extreme eigenvalues.
In an \(M_x\)-orthonormal eigenbasis \(e_j\) of \(\mathcal A_x\),
with real eigenvalues \(\alpha_j\), one obtains
\[
 |F'(x)|
 \leq\sum_j|\alpha_j|\,
       |\langle e_j,(T_x-cI)e_j\rangle_{M_x}|
 \leq\tfrac12d_x\sum_j|\alpha_j|.
\]
This proves the trace-norm bound directly. Let \(L_x\) be the
intersection of the two ranges in SP6, with \(s_x=\dim L_x\geq1\),
and let \(E_{L_x}\) be its orthogonal projection in \(M_x\).
Each of \(\mathcal R_x,\mathcal P_x\) is at least \(E_{L_x}\).
Thus
\(\mathcal A_x=(\mathcal R_x-E_{L_x})-(\mathcal P_x-E_{L_x})\)
is a difference of two positive operators of trace \(2q-s_x\).
The triangle inequality for trace norm gives
\(\|\mathcal A_x\|_1\leq4q-2s_x\leq4q-2\).
Cauchy--Schwarz on its at most \(2q+1\) eigenvalues gives the refinement
\[
 \boxed{
 |F'(x)|\leq\frac{d_x}{2}
 \min\left\{4q-2s_x,\
 \sqrt{(2q+1)
     (8q-4-2\operatorname{Tr}(\mathcal R_x\mathcal P_x))}
     \right\}.}
 \tag{SP10}
\]
The overlap is calculated from the unchanged source and relation
Grams in SP2. Every cross-pairing remains in that expression.

There is a closed integrated bound. Set \(R=M_0^{-1}M^{\rm ar}_{2q}\).
It is positive and self-adjoint in the \(M_0\) inner product, since
\(\langle v,Rv\rangle_{M_0}=\langle v,v\rangle_{M^{\rm ar}_{2q}}>0\).
Let its smallest and largest eigenvalues be \(r_{\min},r_{\max}>0\).
The literal factorization of \(M_x\) gives
\[
 T_x=(I+x(R-I))^{-1}(R-I),\qquad
 d_x=\frac{r_{\max}-1}{1+x(r_{\max}-1)}
           -\frac{r_{\min}-1}{1+x(r_{\min}-1)} .
 \tag{SP11}
\]
For fixed \(x\), the function \((r-1)/(1+x(r-1))\) is increasing
on \(r>0\), its derivative being \((1+x(r-1))^{-2}>0\).
This proves the extreme-eigenvalue formula even though the inner
product varies. Integration, including the identically zero
integrand when an eigenvalue equals one, yields
\[
       \int_0^1d_x\,dx=\log(r_{\max}/r_{\min}).
 \tag{SP12}
\]
Combining SP7--SP12 proves
\[
 \boxed{
 |\Delta_{h,k}|
 \leq\frac12\int_0^1d_x
          \|\mathcal R_x-\mathcal P_x\|_1\,dx
 \leq(2q-1)\log\frac{r_{\max}}{r_{\min}} .}
 \tag{SP13}
\]
The first bound and its overlap refinement SP10 retain more
information than the last bound. When the two Grams are proportional,
\(r_{\max}=r_{\min}\) and the correction is exactly zero, with no
rescaling of their original masses.

SP1--SP13 give finite two-sided control on every original packet.
They do not bound the displayed condition number or projection
overlap along the growing arithmetic family. The next calculation
can use these exact quantities, their original moment formulae,
and the newly retained singular-boundary geometry; none has been
replaced by an assumed asymptotic estimate or an RH verdict.
